\title{Byte Latent Transformer: Patches Scale Better Than Tokens}

\begin{abstract}
This work presents the Byte Latent Transformer ({\textsc BLT}), a novel byte-level LLM architecture that, for the first time, achieves comparable performance to tokenization-based large language models (LLMs) at scale, while offering notable gains in inference speed and robustness. {\textsc BLT} operates by transforming input bytes into adaptively sized patches, which function as the essential computational entities. These patches are delineated by assessing the entropy of forthcoming bytes, thereby allocating greater computational resources and model expressivity in regions exhibiting higher data complexity. We conduct the first scaling study of byte-level models that controls for floating point operations (flops), exploring sizes up to 8B parameters and datasets containing 4T bytes. Our findings affirm that models can be effectively scaled while trained directly on raw bytes, obviating the need for fixed token vocabularies. Dynamic selection of larger patches in predictable contexts yields enhanced efficiency during both training and inference and leads to qualitative advances in complex reasoning and the generalization to rarer patterns. Notably, under fixed inference cost constraints, {\textsc BLT} achieves much improved scaling relative to tokenization-based approaches, by jointly increasing both patch and model capacities.
\end{abstract}

\section{Introduction}
\label{section:intro}

We present the Byte Latent Transformer~(\textbf{BLT}), a novel architecture that eschews tokenization and instead operates directly on sequences of raw bytes. This model achieves, for the first time, performance on par with token-based approaches at large scale, while offering notable gains in both efficiency and resilience (\S\ref{sec:robustness}).

While current large language models (llms) are generally trained in a mostly end-to-end manner, they still rely on tokenization—a heuristic preprocessing stage that segments raw byte sequences into a predetermined set of tokens. This step inherently introduces biases in the compression of strings, giving rise to issues such as sensitivity to domain or modality~\citep{dagangetting}, heightened vulnerability to input noise~(\S\ref{sec:robustness}), insufficient orthographic understanding~\citep{edman2024cute}, and disparities across languages~\citep{liang2023xlm,petrov2024language,limisiewicz2024myte}.

Historically, tokenization has been crucial because training llms directly on raw bytes incurs excessive computational costs at scale, primarily because of the extended sequence lengths involved~\citep{xue2022byt5}. Earlier approaches have attempted to address this challenge by introducing more computationally efficient self-attention mechanisms~\citep{el2020characterbert, clark2022canine} or by utilizing architectures that dispense with attention altogether~\citep{wang2024mambabyte}~(\S\ref{section:related_work}). Nevertheless, such strategies offer significant benefits chiefly for training \textit{smaller models}. When scaling up, the bulk of the computational burden in Transformer architectures is attributable to the expansive feed-forward network layers operating on every byte, rather than to the attention component itself.

To optimize the use of computational resources, we introduce an adaptive, trainable approach that clusters bytes into \textit{patches}~(\S\ref{section:patching}) and design a novel model architecture that integrates both byte-level and patch-level features. In contrast to tokenization, {\textsc BLT} does not rely on a predetermined patch vocabulary. Instead, any collection of bytes can be converted into latent patch representations using efficient, trainable encoder and decoder components. Our findings demonstrate that this strategy enables a \textit{higher} efficiency in computational allocation compared to models that depend on tokenization.

Tokenization-based LLMs distribute identical computational resources to each token, regardless of its individual complexity. This approach favors efficiency at the potential cost of predictive accuracy, as the heuristics used for compression during token induction do not consistently reflect the true difficulty of the prediction tasks. A foundational aspect of our design is the principle that computational effort should be allocated adaptively, concentrating resources where they are most necessary. For instance, generating the last letters of most words is typically straightforward and characterized by lower entropy, making the use of a full-scale transformer unnecessary, in contrast to more challenging cases such as initiating a new sentence. {\textsc BLT}'s architecture~(\S\ref{section:architecture}) embodies this paradigm by employing three transformer modules: two compact byte-level \textit{local models}, complemented by a substantial \textit{latent transformer} that operates globally~(\autoref{fig:arch_high_level}).

To enable the dynamic distribution of computational attention, {\textsc BLT} organizes data by aggregating bytes into patches according to the entropy associated with predicting the next byte. This entropy-aware segmentation allows the model to form context-dependent groupings that maintain a relatively uniform information content across patches.

We conduct the inaugural analysis of flop-efficient scaling for byte-level models, exploring architectures up to 8B parameters trained on 4T bytes. Our results demonstrate that it is feasible to train large-scale models directly on byte sequences, obviating the need for predetermined tokenization schemes. In terms of training flops—measured by total Floating Point OPerations (flops)—{\textsc BLT} achieves parity with Llama 3, while offering up to a 50\% reduction in flops during inference~(\S\ref{section:scaling}).

Moreover, processing raw bytes directly yields marked enhancements in modeling the rarer components of datasets. Models based on {\textsc BLT} exhibit greater resilience to input noise compared to tokenization-based counterparts, as well as improved sensitivity to character-level linguistic features, which we verify through evaluations on orthographic distinctions, phonological tasks, and low-resource machine translation benchmarks~(\S\ref{sec:robustness}).

Additionally, {\textsc BLT} architectures allow for simultaneous increases in both model capacity and patch size without exceeding a constant inference flop limit. Employing longer patch sizes generally leads to computational savings, enabling the allocation of resources toward scaling up the global latent transformer by reducing its frequency of execution. Our controlled scaling experiments with a fixed inference flop budget~(\autoref{fig:fixed_inference_scaling}) reveal that these byte-based models scale more efficiently compared to tokenization-dependent architectures.

To conclude, the primary contributions of this work are as follows: 1) We present {\textsc BLT}, a novel byte latent LLM design capable of dynamic compute allocation to enhance flop efficiency; 2) We provide evidence that our approach reaches flop-equivalent training performance on par with Llama 3 up to the 8B parameter scale, with the added ability to intentionally balance slight decreases in evaluation scores for substantial improvements in flop efficiency—up to 50\%; 3) The {\textsc BLT} architecture offers a new approach for scaling LLMs, permitting increases in model size while keeping the inference budget constant; 4) We empirically validate that {\textsc BLT} exhibits greater resilience to input noise and displays heightened sensitivity to sub-word input details, which token-based LLMs typically overlook. The code for both training and inference with {\textsc BLT} has been made publicly accessible at~\url{https://github.com/facebookresearch/blt}.

\section{Patching: From Individual Bytes to Groups of Bytes}
\label{section:patching}

Dividing bytes into discrete \textit{patches} enables {\textsc BLT} to flexibly distribute computational resources in accordance with contextual needs. Various strategies for partitioning bytes into patches are illustrated in Figure~\ref{fig:patching_types}. More formally, a patching function $f_p$ takes a byte sequence $\pmb{x}=\{x_i,|i=1,\ldots n\}$ of length $n$ and partitions it into $m < n$ patches, denoted $\pmb{p}=\{p_j|j=1,\ldots,m\}$, by assigning each $x_i$ to an indicator set \{0,1\}, where a value of 1 marks the beginning of a new patch. For both token-level and patch-level models, the amount of computation required to handle input data depends chiefly on the number of operations performed by the core Transformer, which, in the case of {\textsc BLT}, is dictated by the number of patches produced by the chosen patching function. As such, the primary determinant of computational expense for processing (during both training and inference) is the mean number of bytes per patch—referred to as the \textit{patch size}—with respect to the applied patching scheme~(\S\ref{section:flops}). In the sections that follow, we present three distinct patching methodologies: fixed-length patching~(\S\ref{section:static-patch}), whitespace-based patching~(\S\ref{section:white-patch}), and adaptive patching driven by entropy estimations from a lightweight byte-level language model~(\S\ref{section:dyn-patch}). We conclude by examining incremental patching and delineating the distinction between tokenization and patching~(\S\ref{section:bpe-patch}).

\subsection{Strided Patching Every K Bytes}
\label{section:static-patch}

A simple method for aggregating bytes is to divide them into patches of a constant size $k$, as implemented in MegaByte~\citep{yu2023megabyte}. This approach, which employs a fixed stride, is straightforward both to train and to use during inference. It also offers an uncomplicated way to adjust the typical patch size, which in turn facilitates precise management of computational costs (flops). Nonetheless, this patching scheme suffers from notable limitations. Primarily, computational resources are not flexibly distributed according to data requirements: for instance, computational steps may be wasted when generating uninformative content like whitespace in code, or insufficient effort may be devoted to bytes conveying substantial information, such as mathematical expressions. Moreover, this strategy can cause inconsistent and context-insensitive segmentations of comparable byte sequences; for example, identical words may end up being fragmented in various ways.

\subsection{Space Patching}
\label{section:white-patch}

\citet{slagle2024spacebyte} introduces an effective modification to strided patching by generating new patches at the occurrence of any space-like byte\footnote{
    Space-like bytes are considered any byte that is not a Latin letter, numeral, or utf-8 continuation byte. Furthermore, at least one byte in each patch must be non-space-like.
}, which frequently correspond to natural linguistic unit boundaries in numerous languages. Within the Space patching approach, each word receives a latent transformer step (i.e., additional computational cost), allowing explicit modeling of every word. This approach ensures consistency in patching words across sequences and focuses computational effort on the challenging predictions that commonly arise after space-like delimiters. For instance, predicting the initial byte of the response to the query ``Who composed the Magic Flute?\uline{\hspace{.6em}}'' is substantially more demanding than anticipating the subsequent bytes after ``M'', because the initial character of the answer greatly narrows the field of probable continuations, thereby making the rest of ``Mozart'' much simpler to generate. Nonetheless, space patching exhibits limitations: it struggles with full generalizability across all languages and domains, and, crucially, it lacks flexibility in patch size adjustment. In the following section, we present a novel patching technique built on the observation that initial bytes of words are generally the most unpredictable, while also allowing for dynamic patch size control.

\subsection{Entropy Patching: Using Next-Byte Entropies from a Small Byte LM}
\label{section:dyn-patch}

Instead of employing rule-based heuristics like whitespace, we adopt a data-driven methodology to detect regions of high uncertainty in next-byte prediction. Our method, termed \textit{entropy patching}, leverages entropy estimations to determine the locations of patch boundaries.

We train a small byte-level auto-regressive language model on the training data for {\textsc BLT} and compute next byte entropies under the LM distribution $p_{e}$ over the byte vocabulary $\mathcal{V}$:

\begin{equation}
H(x_i) = \sum_{v \in \mathcal{V}} p_{e}(x_i=v|\pmb{x}_{<i}) \log p_{e}(x_i=v|\pmb{x}_{<i})
\end{equation}

We explore two approaches for detecting patch boundaries using the entropies $H(x_i)$. The first method selects locations where the entropy surpasses a predefined global threshold, as shown in~\autoref{fig:patching}. The second method focuses on points whose entropy is significantly larger than the preceding value, effectively pinpointing places where the local trend of decreasing entropy within the patch is disrupted.

\begin{align*}
    \text{Global Constraint}\hspace{1cm}H(x_t)& > \theta_g\\
    \text{Approx. Monotonic Constraint}\hspace{1cm}H(x_t)& - H(x_{t-1}) > \theta_r
\end{align*}

Patch boundaries are determined through a simple preprocessing procedure carried out in the dataloading phase. This approach contrasts with the method of~\citet{nawrot-etal-2023-efficient}, where a classifier is trained to detect patch boundaries using entropy predictions. In our experiments~(\S\ref{section:experiments}), we evaluate and contrast these two techniques for classifying bytes as either low or high entropy.

\subsection{The Byte-Pair Encoding (BPE) Tokenizer and Incremental Patching}
\label{section:incremental-patching}
\label{section:bpe-patch}

Numerous contemporary llms, such as our baseline Llama 3, employ subword tokenizers such as bpe~\citep{gage1994new, sennrich-etal-2016-neural}. Throughout this paper, ``tokens'' denote collections of bytes selected from a \textit{finite} set of vocabulary items established before the training process, in contrast to ``patches,'' which describe sequences formed dynamically without reliance on any predetermined vocabulary. A key distinction lies in the fact that, with tokens, the underlying byte-level information is not directly accessible to the model.

A key advancement introduced by {\textsc BLT} compared to traditional tokenization-based approaches lies in its redefinition of the balance between vocabulary size and computational requirements. In conventional llms, expanding the vocabulary typically results in longer average tokens, thereby reducing the number of computational steps needed by the model. However, this also increases the output dimensionality in the model's final projection layer. As a consequence, tokenization-based strategies have limited flexibility in adjusting both token length and the computational demands during inference. As an illustration, Llama 3 increases the average token length from 3.7 to 4.4 bytes, but this comes at the expense of enlarging the embedding table by a factor of four relative to Llama 2.

During generation, {\textsc BLT} must determine at each position in the byte sequence whether a patch boundary has been reached, as this choice dictates if additional computation is required through the Latent Transformer. Crucially, this decision must be made in isolation, without knowledge of bytes that are still to be produced in subsequent steps. Therefore, the method used to segment (or patch) the byte sequence must operate without reference to future bytes. More precisely, we define a patching scheme $f_p$ as *incremental* if it fulfills the following criterion:
$$f_p(\pmb{x}_{<i}) = f_p(\pmb{x})_{<i}$$
bpe does not constitute an incremental patching scheme, since the segmentation of a given prefix may vary depending on the continuation of the sequence, violating the property above\footnote{One can make bpe incremental by introducing a unique delimiter token to specify patch boundaries, though this approach comes at the cost of lengthening the byte sequence.}.

\section{{\textsc BLT} Architecture}
\label{section:architecture}
{\textsc BLT} consists of a primary, large-scale global autoregressive language model utilizing patch-level representations. Accompanying this core model are two compact local models—one responsible for transforming byte sequences into patches, and the other tasked with translating these patch representations back into bytes~(\autoref{fig:arch_high_level}).

\subsection{Latent Global Transformer Model}

\textit{The Latent Global Transformer} refers to an autoregressive transformer architecture, denoted as $\mathcal{G}$ and composed of $l_{\mathcal{G}}$ layers. It transforms a sequence of latent input patch representations, $p_j$, into corresponding output patch representations, $o_j$. In this work, the notation uses $j$ for patches and $i$ for bytes. The model utilizes a block-causal attention mask~\citep{dubey2024llama}, confining the attention mechanism to the current patch and its predecessors within the current document. The computational cost of pre-training and inference is primarily attributed to this model; therefore, selectively activating $\mathcal{G}$ enables fine-grained management of computational resources across various segments of the input and output depending on their complexity.

\subsection{Local Encoder}
\label{section:local}

\textit{The Local Encoder Model}, represented as $\mathcal{E}$, is a compact transformer variant consisting of $l_{\mathcal{E}} << l_{\mathcal{G}}$ layers, specifically designed for the rapid transformation of a byte input sequence $b_i$ into informative patch representations $p_j$. Unlike standard transformers, $\mathcal{E}$ integrates a cross-attention mechanism after each transformer layer, which functions to aggregate the byte-level embeddings into the higher-level patch representations (see~\autoref{fig:crossattn}).

Initially, the bytes $b_i$ are projected into embedding vectors $x_i$ via a lookup matrix of shape $\mathbb{R}^{256 \times h_{\mathcal{E}}}$. These base embeddings may be further supplemented by hash-embeddings for incorporating extra contextual data~(\S\ref{sec:hash-emb}).

Subsequent layers—alternating between standard transformer blocks and cross-attention mechanisms~(\S\ref{sec:enc-cross-attn})—incrementally convert the sequence into patch representations $p_i$. These patch representations are subsequently consumed by the global transformer, $\mathcal{G}$.

For the attention scheme, each transformer layer is governed by a \textit{local block causal} mask that restricts each byte to attend only to a window of the $w_{\mathcal{E}}$ most recent prior bytes. While this window can span across dynamic patch borders, it is not permitted to extend beyond document edges. The following subsections provide further elaboration on the embedding process and the architecture of the cross-attention modules.

\subsubsection{Encoder Hash n-gram Embeddings}
\label{sec:ngram-emb}
\label{sec:hash-emb}
To develop resilient and informative representations for each position $i$, it is essential to leverage contextual information from earlier bytes. Within {\textsc BLT}, this is realized by encoding both the individual byte $b_i$ and its contribution to a byte-level n-gram structure. Specifically, at each step $i$, we begin by assembling byte-grams

\begin{align}
    g_{i,n}=\{b_{i-n + 1},\ldots, b_i\}
\end{align}

for each byte index $i$, considering $n$ ranging from three to eight.\footnote{
If $i < n$, we do not include byte-grams of length $n$ or greater.
}

Subsequently, we propose hash $n$-gram embeddings, wherein all byte $n$-grams are transformed through a hash function into an index within a corresponding embedding table $E_{n}^{hash}$ of fixed dimensionality, for each $n \in \{3,4,5,6,7,8\}$~\citep{bai2010learning}. The produced embedding is combined with the original byte embedding; the resultant sum is then normalized and provided as input to the local encoder. The combined embedding is computed as follows:

\begin{align}
e_i &= x_i + \sum_{n = 3,...,8}E_{n}^{hash}(\text{Hash}(g_{i,n})) \\
\text{where, Hash}(g_{i,n}) &= \text{RollPolyHash}(g_{i,n}) \% |E_{n}^{hash}|
\end{align}

To compute $e_i$, we divide by the total number of $n$-gram sizes incremented by one, utilizing RollPolyHash as outlined in Appendix~\ref{appendix:rollpolyhash}. Section~\ref{section:ablations} presents ablation studies examining how varying the $n$-gram length and the embedding table size in $n$-gram hash embeddings impacts the scaling law behavior under fixed computational budgets. Beyond hash-based $n$-gram embeddings, we further investigated embeddings constructed from $n$-gram frequencies; a comprehensive discussion of these experiments is provided in Appendix~\ref{appendix:ngrams}.

\subsubsection{Encoder Multi-Headed Cross-Attention}
\label{sec:enc-cross-attn}
Our approach closely parallels the input cross-attention component from the Perceiver architecture~\citep{jaegle2021perceiver}, with the key distinction that here, latent representations are aligned with dynamic patch representations rather than being tied to a static collection of latents~(\autoref{fig:crossattn}). Each latent attends exclusively to the bytes contained within its corresponding patch. Specifically, for every patch $p_j$, a query vector is created by aggregating the byte representations within $p_j$—for example, through a pooling operation—then applying a linear transformation defined by $\mathcal{E}_{C} \in \mathbb{R}^{h_{\mathcal{E}} \times (h_{\mathcal{E}} \times U_{\mathcal{E}})}$, where $U_{\mathcal{E}}$ denotes the number of cross-attention heads used in the encoder. In this setup, if $f_{\text{bytes}}(p_j)$ signifies the ordered set of bytes pertaining to patch $p_j$, the following computation is performed:

\begin{align}
P_{0,j} &= \mathcal{E}_{C}(f_\text{bytes}((p_j)), f~ \text{is a pooling function} \\
P_l &= P_{l-1} + W_o\left(\text{softmax}\left(\frac{QK^T}{\sqrt{d_k}}\right)V\right) \\
\text{where } Q_j &= W_q(P_{l-1,j}), K_i = W_k(h_{l-1,i}), V_i = W_v(h_{l-1,i}) \\
h_l &= \text{Encoder-Transformer-Layer}_l(h_{l-1})
\end{align}

Here, $P \in \mathbb{R}^{n_p \times h_{\mathcal{G}}}$ denotes the representations of $n_p$ patches that serve as input to the global model. These representations are initialized by aggregating the byte embeddings $e_i$ associated with each patch $p_j$. The matrices $W_q$, $W_k$, $W_v$, and $W_o$ act as the linear transformations for the queries, keys, values, and output, where keys and values are obtained by projecting the byte-level representations $h_i$ from the previous layer (using $e_i$ in the initial layer). To ensure that queries focus only on relevant information, we implement a masking technique tailored to patches such that a query $Q_j$ attends exclusively to the keys and values linked to the bytes in patch $j$. Since multi-headed attention is applied over $Q$, $K$, and $V$, and patch representations often have higher dimensionality ($h_{\mathcal{G}}$) than $h_{\mathcal{E}}$, $P_l$ is structured as multiple attention heads, each with dimension $h_{\mathcal{E}}$ during cross-attention; these heads are subsequently concatenated to reconstruct the $h_{\mathcal{G}}$-dimensional output. Furthermore, we apply a pre-LayerNorm to queries, keys, and values, and omit any positional embeddings within this cross-attention layer. The architecture includes a residual connection that wraps around the entire cross-attention component.

\subsection{Local Decoder} 

Much like the local encoder, the local decoder $\mathcal{D}$ also utilizes a compact transformer architecture, comprising $l_{\mathcal{D}} << l_{\mathcal{G}}$ layers. This decoder translates a sequence of global patch representations $o_j$ into the original sequence of bytes, $y_i$. The local decoder generates each byte sequentially, conditioned on the bytes decoded so far. Therefore, its input consists of the hidden states produced by the local encoder for the given byte sequence. Within the decoder, $l_{\mathcal{D}}$ layers are structured to alternate between cross-attention modules and standard transformer layers. The cross-attention module precedes the transformer layer, facilitating the mapping from patch-level representations to byte-level representations, after which the transformer processes this byte-level sequence.

\subsubsection{Decoder Multi-headed Cross-Attention} 

In the decoder cross-attention, the roles of the queries and key/values are interchanged i.e. the byte-representations are now the queries, and the patch representations are now the key/values. The initial byte-representations for the cross-attention are initialized as the byte embeddings from the last encoder layer i.e. $h_{l_{\mathcal{E}}}$. The subsequent byte-representations for layer $l$, $d_{l,i}$ are computed as:

\begin{align}
D_0 &= h_{l_{\mathcal{E}}} \\
B_l &= D_{l-1} + W_o\left(\text{softmax}\left(\frac{QK^T}{\sqrt{d_k}}\right)V\right), \\
  &\text{where } Q_i = W_q(d_{l-1,i}), K_i = W_k(\mathcal{D}_C(o_j)), V_i = W_v(\mathcal{D}_C(o_j)) \\
  D_l &= \text{Decoder-Transformer-layer}_l(B_l)
\end{align}

Here, $W_k$ and $W_v$ once more serve as the key and value projection matrices, respectively, both implementing linear transformations in conjunction with the split operation $\mathcal{D}_C$ on the ultimate patch outputs $o_j$ obtained from the global model. The query projection matrix, $W_q$, transforms byte representations $d_{l-1}$—originating from the previous decoder transformer layer (or $h_{l_{\mathcal{E}}}$ when considering the initial layer). The output projection is carried out by $W_o$. As a result, we have $B \in \mathbb{R}^{h_{\mathcal{D}} \times n_b}$, where $n_b$ denotes the total number of output bytes. The subsequent layer of decoder representations, $D_l$, is generated by feeding the cross-attention module’s output $B$ into a decoder transformer layer. Consistent with our approach in the local encoder cross-attention, we deploy multiple attention heads, apply pre LayerNorms, omit positional embeddings, and maintain a residual connection across the cross-attention module.

\section{Experimental Setup}
\label{section:experiments}
We have meticulously constructed controlled experiments to rigorously compare {\textsc BLT} with models using traditional tokenization, ensuring that {\textsc BLT} is not afforded any undue advantages through access to extended sequence lengths.

\subsection{Pre-training Datasets} In our experiments, all model variants are pre-trained on two principal datasets: 1) The Llama 2 dataset~\citep{touvron2023llama}, which contains 2 trillion tokens amassed from diverse publicly available data sources that have undergone extensive cleaning and filtering processes to enhance overall quality; and 2) {\textsc BLT}-1T: A novel dataset consisting of 1 trillion tokens sourced from an assortment of public datasets, which additionally incorporates a selected portion of the Datacomp-LM pre-training dataset~\citep{li2024datacomp}. The Llama 2 dataset serves as the basis for scaling law analyses pertaining to the optimal token count, following methodology in~\cite{dubey2024llama}, to inform architectural decisions in {\textsc BLT}, whereas {\textsc BLT}-1T is utilized for comprehensive pre-training in order to directly benchmark performance with Llama 3 on downstream evaluation tasks. It is important to note that neither dataset includes information derived from Meta’s proprietary products or services. Moreover, in our baseline studies involving tokenizer-based architectures, we employ the Llama 3 tokenizer with a 128K vocabulary, as this demonstrated superior baseline results compared to the Llama 2 tokenizer across our empirical analyses.

\subsection{Entropy Model} For our experiments, the entropy model is implemented as a byte-level language model, trained on the identical dataset as the comprehensive {\textsc BLT} model. By default, the architecture utilized is a transformer composed of 100M parameters, organized into 14 layers, featuring a hidden size of 512, and employing sliding window attention spanning 512 bytes. Other hyperparameter settings are aligned with those used in our local and global transformer baselines. Additional experiments, detailed in \autoref{section:ablations}, involved varying the model size, receptive field, and architectural choices. Notably, if the model's receptive field is reduced sufficiently, it becomes possible to represent the trained entropy model as a compact and efficient lookup table.

\subsection{Entropy Threshold and Equalizing Context Length} For models employing entropy-based patching, we determine a patching threshold that yields a target average \textit{patch size} on the pretraining dataset mixture. In {\textsc BLT}, as opposed to tokenization, the selection of \textit{patch size} is flexible, directly affecting the effective context length utilized by the model. To ensure fairness and prevent models with larger patch sizes from benefiting from increased context, we enforce an equivalent average byte count per batch across models. Consequently, when a larger patch size is used, we proportionally decrease the number of sequences per batch. Specifically, for Llama 2 data, an 8k byte context length is adopted, whereas for the {\textsc BLT}-1T dataset, the context length is set to 16k bytes on average; in both cases, the batch size is held constant at an average of 16M bytes.

Although the average batch size remains fixed, the proportion of bytes to patches varies across batches when using dynamic patching approaches. To improve computational efficiency, our implementation of {\textsc BLT} organizes batches so that they are filled with a consistent number of patches, thereby eliminating the need for padding within the more computationally intensive latent transformer. This approach guarantees that each batch maintains an equal count of patches. Throughout training, byte sequences are padded or occasionally truncated—set to 12k bytes for Llama 2 data and 24k bytes for the {\textsc BLT}-1T dataset—in order to prevent memory usage spikes caused by sequences containing particularly large patches.

\subsection{Entropy Model Context}
\label{section:entropy-context}
Our observations indicate that entropy patching leads to incrementally larger patches in structured formats such as multiple choice tasks (refer to the patching applied in an MMLU example in \autoref{fig:mmlu_patching}), where repetitiveness is common. These patch size differences stem from the entropy model context displaying reduced entropy for recurring content. To address this in our large-scale implementation of {\textsc BLT}-Entropy with a patch size of 4.5, we refresh the entropy context by introducing new line breaks and implement an approximate monotonicity constraint, as this method is less affected by the “entropy drift” resulting from varying context lengths. Notably, this adjustment influences only the entropy calculation process; the method for determining the entropy threshold remains unchanged.

\subsection{FLOPs Estimation} 
\label{section:flops}

Our methodology primarily adopts the transformer floating point operation (flop) calculations outlined in Chinchilla~\citep{hoffmann2022training}, which account for flops consumed by the feed-forward networks, the qkvo projections within the self-attention mechanism, as well as the attention calculation and subsequent output projection. In contrast to the referenced work, however, we posit that the input embedding layer utilizes an optimized lookup table in place of a conventional dense matrix multiplication, effectively rendering its flop contribution negligible. Consistent with established practice, we approximate the backward pass as incurring double the flops of the forward pass.

To compute flops \textit{per byte} for {\textsc BLT} models, we add up the flops for the local encoder transformer, the global latent transformer, and the local decoder transformer, together with the cross attention blocks in the encoder and the decoder:

\begin{align}
    \text{FL}_{\text{{\textsc BLT}}} &= \text{Transf. FL}(h_{\mathcal{G}}, l_{\mathcal{G}}, m=n_{ctx}/n_p,V=0) / n_p\\
    &\quad +\text{Transf. FL}(h_{\mathcal{E}}, l_{\mathcal{E}}, m=w_{\mathcal{E}}, V=0) \\
    &\quad + \text{Transf. FL}(h_{\mathcal{D}}, l_{\mathcal{D}}, m=w_{\mathcal{D}}, V=256) \\
    &\quad + \text{Cross Attn. FL}(h_{\mathcal{E}}, l_{\mathcal{E}}, m=n_p, r=n_p/k) \cdot (k/n_p) \\
    &\quad + \text{Cross Attn. FL}(h_{\mathcal{D}}, l_{\mathcal{D}}, m=k, r=k/n_p)
\end{align}

Here, $n_{ctx}$ represents the sequence length in bytes, and $n_p$ denotes the patch size. The parameter $r$ specifies the proportion of queries relative to key/values, while $k$ indicates the ratio of patch dimension to byte dimension; that is, $k$ reflects the count of local model partitions combined to yield a complete model representation ($k=2$ in \autoref{fig:crossattn}). The variable $V$ stands for the vocabulary size used during output projection, which is incorporated exclusively within the local decoder. The context length for attention, $m$, varies based on whether the module is acting upon the byte sequence or the patch sequence. To account for this, the computation of attention flops is adjusted for each architecture component. The precise formulae for calculating Transformer-FLOPs and Cross-Attention FLOPs are given in Appendix~\ref{section:flops-appendix}.

\subsection{Bits-Per-Byte Estimation} Perplexity only makes sense in the context of a fixed tokenizer as it is a measure of the uncertainty for each token. When comparing byte and token-level models, following previous work~\citep{xue2022byt5, yu2023megabyte, wang2024mambabyte}, we instead report Bits-Per-Byte (BPB), a tokenizer independent version of perplexity. Specifically:

\begin{align}
    \text{BPB}(x)&= \frac{\mathcal{L}_{CE}(\pmb{x})}{\ln(2)\cdot n_{\text{bytes}}}
\end{align}

Here, the cumulative cross-entropy loss representing the uncertainty of the data $\pmb{x}$ is divided by both a constant factor and the overall number of bytes in $\pmb{x}$ for normalization purposes.

\subsection{Transformer Architecture Hyperparameters} In configuring all transformer components within {\textsc BLT}—encompassing both local and global modules—we adopt most architectural conventions from Llama~3~\citep{dubey2024llama}. The feed-forward sublayers incorporate the SwiGLU activation function~\citep{swiglu}, while self-attention submodules employ rotary positional embeddings (RoPE)~\citep{RoPe} set with $\theta=500000$~\citep{xiong2024effective}. For normalization, RMSNorm~\citep{RMSNorm} is consistently utilized. All self-attention modules that leverage predetermined attention patterns—such as \textit{block causal} and \textit{fixed-window block causal}—are implemented using Flash attention~\citep{dao2022flashattention}, with attention windows fixed at 512 for the latter. Given that the cross-attention layers require patch-specific, dynamic masking, Flex Attention\footnote{\url{https://pytorch.org/blog/flexattention}} is employed for its fused execution, thereby substantially accelerating the training process.

\subsection{{\textsc BLT}-Specific Hyperparameters}

To evaluate the performance of {\textsc BLT} models, our experimental framework investigates two primary aspects: scaling behavior and downstream task performance, considering model capacities at 400M, 1B, 2B, 4B, and 8B parameters. Details of the model architecture hyperparameters can be found in Appendix~Table~\ref{tab:arch}. For initializing queries in the initial cross-attention module within the local encoder, max-pooling is employed. Across all {\textsc BLT} variants, $500,000$ hashes are used with a single hash function, and n-gram sizes span from 3 to 8. The learning rate is set uniformly at $4e-4$ for all models. Selection of this learning rate is informed by a hyperparameter sweep conducted at the 400M and 1B model scales, ranging between $1e-3$ and $1e-4$, indicating that identical learning rates are optimal for both the token and {\textsc BLT} models. For scaling trend analyses using Llama-2 datasets, we adopt batch sizes either in accordance with the guidelines of~\cite{dubey2024llama} or in terms of equivalent byte counts. The AdamW optimizer~\citep{adamw} is utilized for model optimization, configured with $\beta_1 = 0.9$, $\beta_2 = 0.95$, and $\epsilon=10^{-8}$. A linear warm-up over 2000 steps precedes cosine learning rate decay to zero. Additionally, a weight decay factor of 0.1 is applied, with global gradient norm clipping enforced at 1.0.

\section{Scaling Trends}
\label{section:scaling}

We offer a comprehensive analysis of how byte-level models scale, providing guidance for the continued scaling of {\textsc BLT} models. Our investigation into scaling seeks to overcome several shortcomings of earlier studies on byte-level modeling: (a) We examine performance in compute-optimal training conditions, (b) We develop parallel 8B models using substantial datasets—reaching up to 1T tokens or 4T bytes—and assess them on downstream benchmarks, and (c) We explore scaling behavior under constraints where inference costs are held fixed. Further on, we will explore the unique benefits conferred by processing byte-sequences directly.

\subsection{Parameter Matched Compute Optimal Scaling Trends}
\label{section:parameter-matched-scaling}
We conduct experiments on the Llama 2 dataset by training several \textit{compute-optimal} bpe and {\textsc BLT} models at four different scales, spanning from 1B to 8B parameters. Training flops are plotted alongside language modeling accuracy on a selected subset of the mixed training data. In alignment with the protocol defined by Llama~3~\citep{dubey2024llama}, bpe models are trained using the optimal proportion between model size and training corpus. This approach, designed to be \textit{compute-optimal}, is expected to yield maximum attainable performance for a specified computational expenditure~\citep{hoffmann2022training}, thereby serving as a strong reference point. For every bpe model, a corresponding {\textsc BLT} model is also trained on identical data, employing a Latent Transformer with architecture and size closely matching that of its bpe Transformer counterpart.

As shown in~\autoref{fig:scaling-compute-opt} (right), {\textsc BLT} models consistently achieve equal or superior performance relative to their bpe-based equivalents, and this pattern persists as model size and computational requirements increase. According to our knowledge, {\textsc BLT} represents the first byte-level Transformer design to exhibit scaling behavior on par with BPE models under compute-optimal settings. This finding supports our hypothesis that the ideal proportion between parameters and training compute for bpe also extends to {\textsc BLT}, or at a minimum, does not diverge significantly.

Achieving competitive bpe scaling requires both advancements in model architecture and the implementation of dynamic patching strategies. In ~\autoref{fig:scaling-compute-opt} (left), we present a comparison between space-patching models and Llama 3. For the analysis, we estimate the performance of SpaceByte \citep{slagle2024spacebyte} by employing {\textsc BLT} space-patching, but exclude n-gram embeddings and cross-attention mechanisms. While SpaceByte exhibits gains relative to Megabyte, its performance is still substantially below that of Llama 3. \autoref{fig:scaling-compute-opt} (right) demonstrates the impact of both architectural refinement and dynamic patching, revealing that, at larger scales, {\textsc BLT} models achieve results comparable to leading tokenizer-based systems like Llama 3.

Additionally, our results indicate that for models relying on tokenization, selecting the Llama-3 tokenizer leads to better performance compared to employing the Llama-2 tokenizer when trained on identical datasets.

In summary, the {\textsc BLT} architecture demonstrates a scaling pattern that falls between those observed for Llama 2 and Llama 3 when adopting notably larger patch sizes. While Llama 2 and 3 utilize bpe tokenizers with mean token lengths of 3.7 and 4.4 bytes, respectively, {\textsc BLT} is able to match these scaling behaviors with mean patch sizes reaching up to 6 or 8 bytes. Since inference flop decreases as the average patch size increases, selecting an 8-byte patch size can reduce inference computational cost by approximately 50\%. Moreover, models with increased patch sizes tend to yield improved results as both model and dataset sizes are enlarged. When using 8-byte patches, {\textsc BLT} initially underperforms relative to bpe Llama 2 at a model scale of 1B parameters, but it surpasses bpe performance at the 7B scale. These outcomes imply that larger patch sizes could be even more advantageous at higher scales, indicating the potential feasibility of adopting even greater patch sizes as model capacity and available computation expand.

\subsection{Beyond Compute Optimal Task Evaluations}
\label{section:evals}
To further investigate scaling behaviors, we extend the training of an 8B {\textsc BLT} model past the compute-optimal regime, utilizing the larger, higher-quality {\textsc BLT}-1T dataset. We then assess its capabilities across a range of well-established benchmarks that test both classification and generation abilities. For this evaluation, we focus on tasks measuring common sense reasoning, world knowledge, and code generation:

\paragraph{Classification tasks} cover ARC-Easy (0-shot)~\citep{Eval_ARC}, Arc-Challenge (0-shot)~\citep{Eval_ARC}, HellaSwag (0-shot)~\citep{Eval_hellaswag}, PIQA (0-shot)~\citep{Eval_piqa}, and MMLU (5-shot)~\citep{hendrycks2020measuring}. Here, we use a prompt-scoring approach, where we compute the likelihood for each possible answer and report average accuracy.

\paragraph{Coding related generation tasks:} Performance on code generation is measured by pass@1 results for MBPP (3-shot)~\citep{Eval_MBPP} and HumanEval (0-shot)~\citep{Eval_HumanEval}, providing an estimate of the model’s proficiency in generating correct Python code.

In \autoref{tab:evals}, we present a comparison among three different models trained on the {\textsc BLT}-1T dataset: one employing the Llama 3 tokenizer with byte pair encoding,\footnote{We opt for the Llama 3 tokenizer, given its 128k vocabulary outperforms that of Llama 2, which uses a 32k vocabulary.} as well as two {\textsc BLT} model variants. The first variant applies a space-patching approach ({\textsc BLT}-Space), while the second leverages an entropy-based patching strategy ({\textsc BLT}-Entropy), imposing an approximate monotonicity constraint and reinitializing the entropy model’s context whenever a newline is encountered, as outlined in~\autoref{section:entropy-context}. All three models were allotted comparable flop budgets during training. Notably, for {\textsc BLT}-Entropy, we also alter the entropy threshold during inference, reducing it from 0.6 to 0.1, which leads to better task performance though it necessitates an increased number of inference steps.

The {\textsc BLT}-Entropy model surpasses Llama 3 on 4 of the 7 evaluated tasks, despite both models being trained with an equal quantity of bytes. This performance gain can likely be attributed to two primary factors: (1) more efficient utilization of training compute through dynamic patching, and (2) direct byte-level representation, rather than token-based modeling.

Conversely, {\textsc BLT}-Space demonstrates lower performance than the Llama 3 tokenizer on nearly every benchmark except one. However, due to its substantially larger mean patch size of 6 bytes, it is able to notably decrease inference flops. For context, models utilizing bpe and entropy-patching approaches exhibit similar mean patch sizes of about 4.5 bytes on the training dataset mixture. Given an identical training budget, the model employing a larger patch size processes 30\% more data than the other two, which may drive {\textsc BLT} to deviate further from the compute-optimal regime.

\subsection{Patches Scale Better Than Tokens} In {\textsc BLT} models, it is possible to scale both the model architecture and the patch size concurrently, while holding constant both the overall computational costs for training and inference, as well as the volume of training data. The flexibility to arbitrarily enlarge patch size distinguishes patch-based models, as they are not bound by the efficiency limitations characteristic of fixed-vocabulary token-based approaches, as explained in Section~\ref{section:bpe-patch}. Increasing the length of patches reduces computational overhead, thus allowing for additional resources to be allocated to expanding the global latent transformer, since its invocation frequency diminishes as patch size grows.

A controlled scaling analysis is performed to evaluate whether, at a constant inference budget, increasing model size while reducing the number of inference steps on larger data patches yields better performance compared to using smaller models with more steps. Using baseline models of 400m and 3.6B parameters equipped with the Llama 2 tokenizer, we identify flop-equivalent configurations utilizing the Llama 3 tokenizer, as well as {\textsc BLT}-Entropy models averaging patch sizes of 6 and 8 bytes within the same training data mix (refer to~\autoref{tab:fixed-inf-params} for specific model configurations). For models with a patch size of 8, the number of encoder layers is increased to 3 as opposed to 1. Each configuration is trained using several training flop constraints.

As illustrated in \autoref{fig:fixed_inference_scaling}, {\textsc BLT} models exhibit more favorable scaling behavior compared to tokenization-based models across both inference flop classes. Initially, BPE models outperform at lower training budgets but are quickly overtaken by {\textsc BLT} models soon after the compute-optimal point is surpassed. In real-world scenarios, allocating greater resources towards pretraining can be advantageous, yielding higher quality models under a fixed inference constraint. A salient illustration is the category of 8B models, such as Llama 3.1, which underwent training on data volumes nearly two orders of magnitude beyond what is considered compute-optimal for its scale.

The point at which {\textsc BLT} surpasses token-based models in performance now occurs at a ratio slightly nearer to the compute-optimal threshold for larger flop class models; specifically, the crossover shifts from approximately 3 times to 2.5 times the optimal compute allocation. In addition, for models employing a patch size of 8, the scaling advantage becomes more pronounced with increased computational capacity, enabling these models to outperform others more rapidly. As highlighted in Section~\ref{section:parameter-matched-scaling}, scaling models with larger patch sizes generally narrows the performance gap with BPE-based models at higher parameter counts. One contributing factor is the declining proportion of overall flops dedicated to the byte-level Encoder and Decoder components, which exhibit slower scaling compared to the Latent Transformer. Notably, a twentyfold increase in model parameters, from 400M to 8B, results in only a twofold increase in the local parameters associated with {\textsc BLT}. This distinction is notable because adjustments in patch size influence only the computational requirements of the Latent Transformer, leaving byte-level components unaffected. This dynamic helps explain why the relative size of {\textsc BLT}-Entropy with ps=8 increased modestly from 1.6x to 1.7x compared to Llama 2 as model scale was enlarged.

To summarize, our investigation into the scaling of patch length reveals that the {\textsc BLT} architecture, which leverages a patch-based framework, attains enhanced scaling behavior when patch size and model capacity are jointly increased. Notably, these scaling improvements continue—and may become more pronounced—as the model size grows.

\section{Byte Modeling Improves Robustness}
We further assess the robustness of the {\textsc BLT} model relative to traditional token-based counterparts that do not have direct access to byte-level features. Additionally, we propose a methodology to transform pretrained token-based models into byte-compatible forms.
\label{sec:robustness}

\subsection{Character-Level Tasks}

Initial efforts to develop byte-level models were largely driven by their inherent resistance to noise at the byte level in the input data, as well as their capacity to recognize and utilize sub-token structures—an area where models reliant on tokenizers often face limitations. To empirically assess these characteristics, we conduct supplementary experiments on various benchmarks designed to test model robustness against noisy inputs and evaluate sensitivity to character-level features, spanning English and multilingual characters, as well as digits and phonemes. The outcomes of these evaluations are summarized in Table~\ref{tab:char_tasks}.

\paragraph{Noisy Data} To assess robustness, we introduce noise into the benchmark classification datasets outlined in Section~\ref{section:evals}, generating altered versions for comparison between tokenizer-based approaches and {\textsc BLT}. Five distinct character-level noise schemes are utilized: (a) \textit{AntSpeak}: the input text is converted fully to uppercase and each character is separated by a space; (b) \textit{Drop}: 10\% of the characters within the text are randomly deleted; (c) \textit{RandomCase}: each character in the text is randomly switched to either uppercase or lowercase with a 50\% probability; (d) \textit{Repeat}: a selection of 20\% of the characters are repeated up to four times; (e) \textit{UpperCase}: every character in the text is capitalized. For evaluation, each noise type is separately applied to either the prompt, the completion, or both, and we record mean performance across these scenarios. Table~\ref{tab:char_tasks} presents outcomes for noisy HellaSwag~\citep{Eval_hellaswag}, demonstrating that {\textsc BLT} consistently shows greater robustness compared to tokenizer-based baselines, achieving an average margin of 8 points over a model trained on identical data, and also surpasses the Llama 3.1 model, which was trained on a significantly larger corpus.

\paragraph{Phonology - Grapheme-to-Phoneme (G2P)} We evaluate {\textsc BLT}'s proficiency in converting a sequence of graphemes (characters of a given word) to its corresponding phonemic transcription (pronunciation in phonemes). Table \ref{tab:char_tasks} displays performance metrics for the G2P task under a 5-shot learning scenario with Phonology Bench~\citep{suvarna2024phonologybench}. The findings indicate that {\textsc BLT} achieves superior results compared to the Llama 3 1T tokenizer-based baseline on this task.

\paragraph{CUTE}
To evaluate the ability of {\textsc BLT} to process information at the character level, we conduct experiments using the CUTE benchmark~\citep{edman2024cute}. This benchmark consists of multiple tasks grouped into three main types: compositional understanding, orthographic similarity recognition, and sequence manipulation capabilities. These tasks prove particularly challenging for standard tokenizer-based models, which tend to know how their internal tokens are spelled but generally lack the capacity to leverage such information for effective text manipulation. As presented in Table~\ref{tab:char_tasks}, {\textsc BLT}-Entropy surpasses both BPE Llama 3 variants by over 25 points on the CUTE benchmark. Notably, our model shows outstanding aptitude in character manipulation, attaining 99.9\% accuracy in both spelling-related tasks. This level of improvement is achieved despite {\textsc BLT} being trained on only 1/16th of the data used for Llama 3.1, highlighting that tokenization methods like BPE are less adept at learning character-level patterns. Figure \ref{fig:cute_examples} provides examples that reveal situations where the Llama 3 tokenizer model falls short, while the {\textsc BLT} model handles them effectively. It is worth noting that word deletion and word insertion are the only instances where the BPE approach exhibits superior results; although these word-level manipulations may be inherently simpler for a BPE system, the performance difference is marginal, suggesting that constructing word-level representations from characters is likely easier than extracting character-level patterns from preexisting words. All evaluations follow an identical experimental protocol, utilizing the original Huggingface prompts. It is possible that further gains could be realized for BPE models with improved prompt engineering.

\paragraph{Low Resource Machine Translation} We assess the performance of {\textsc BLT} for both directions of translation involving six widely spoken language families and twenty-one under-resourced languages—spanning different writing systems—using the FLORES-101 benchmark~\citep{goyal2022flores}. SentencePiece BLEU scores are presented in Table~\ref{tab:flores}. The findings indicate that {\textsc BLT} consistently exceeds the accuracy of a model employing the Llama 3 tokenizer, achieving a 2-point improvement in translations into English and a 0.5-point enhancement when translating from English. For frequently translated language pairs, {\textsc BLT} matches or slightly surpasses the results of Llama 3. Importantly, for a broad range of language pairs within low-resource families, {\textsc BLT} yields superior outcomes relative to Llama 3, highlighting the utility of byte-based modeling in better generalizing to languages with rare or atypical byte sequences.

\subsection{Training {\textsc BLT} from Llama 3}
\label{sec:distillation}
We investigate a strategy wherein {\textsc BLT} models benefit from pre-existing pre-trained tokenizer-based architectures by initializing the global transformer parameters with those taken from a pre-trained Llama 3.1 model. This initialization facilitates faster convergence and enhanced training efficiency. In the subsequent training phase, the global transformer's parameters are fine-tuned using a learning rate set to one-tenth that used for updating the local encoder and local decoder, and training proceeds for the number of steps optimal for Llama 3. Performance outcomes, as reported in Table~\ref{tab:distill}, indicate that {\textsc BLT} initialized from Llama 3.1 yields consistently better results compared to both the standard Llama 3 and a {\textsc BLT} baseline, both of which consume an equivalent amount of computational resources. Furthermore, even against our {\textsc BLT}-Entropy variant trained on a much larger dataset (1T tokens, see Table~\ref{tab:evals}), {\textsc BLT} initialized from Llama 3.1 demonstrates enhanced results on the MMLU benchmark. This highlights the method's potential for drastically reducing training compute requirements while achieving superior task performance.

This approach can be interpreted as converting models that rely on tokenizers into ones that do not, thereby adapting a pre-trained LLaMA 3.1 model into a {\textsc BLT} architecture. For a thorough evaluation, we present results for the original LLaMA 3.1, which was trained on 15T tokens, alongside those for {\textsc BLT} obtained from LLaMA 3, as shown in Table \ref{tab:distill}. While our model exhibits only a modest decrease in accuracy on MMLU and HumanEval, there is a more pronounced decline on several other benchmarks. These findings indicate the need for future research focused on enhancing the benefits of the pre-trained model, with an emphasis on refining data mixture strategies and optimizing additional hyperparameters.

\section{Ablations and Discussion}
\label{section:ablations}
Here, we present a series of ablation studies aimed at motivating the architectural decisions made for {\textsc BLT}, as well as examining the patching mechanism and the selection of hyper-parameters used in the {\textsc BLT} 8B-parameter model trained on the {\textsc BLT}-1T dataset.

\paragraph{Entropy Model Hyper-parameters} To investigate how the size of the entropy model and the length of the context window affect scaling behavior, we conduct experiments using byte-level entropy transformer architectures. These models vary in scale from 1m up to 100m parameters and are trained with context windows ranging from 64 to 512 tokens. We generate and present scaling law plots of bits per byte (bpb) against training FLOPs for our $400m$ and $1b$ {\textsc BLT} models, utilizing the Llama-2 dataset, as shown in \autoref{fig:entropy_ablation}. Our results indicate that increasing both model size and context length generally leads to improved scaling performance; however, the gains become marginal when model sizes exceed 50m parameters.

\paragraph{Types of Patching}

We conduct an ablation study on the four patching strategies outlined in Section~\ref{section:patching}: 1) Strided Patching using strides of 4 and 6, 2) Whitespace-based Patching, 3) BPE Tokenizer patching utilizing the Llama 3 tokenizer, and 4) Entropy-guided patching employing a compact byte-level language model.

Although dynamic patching effectively shortens sequence lengths, we ensure that the context length remains comparable across all patching strategies by standardizing sequence length. Consequently, each model is exposed to an equivalent number of bytes per sequence throughout both training and inference, thus eliminating confounding variables related to context size. The outcomes of these controlled experiments are illustrated in Figure~\ref{fig:scaling-compute-opt}. Notably, every patching scheme other than static patching demonstrates superior performance, with space patching yielding results nearly indistinguishable from dynamic entropy-based patching.

\autoref{tab:evals_ablation} summarizes benchmark results for {\textsc BLT} models, comparing the performance of tokenizer-based approaches, space patching, and entropy-based patching, all trained on the Llama 2 dataset for an empirically determined number of steps~\citep{dubey2024llama}. While space patching employs a more straightforward method that circumvents the need for real-time entropy model computations during training, our experiments confirm that the improvements yielded by entropy-based patching on scaling metrics (see Section~\ref{section:scaling}) are also reflected in performance across downstream benchmarks.\footnote{The space patching data are derived from earlier training runs that did not utilize cross-attention; however, consistent patterns persist in experiments where cross-attention is employed.}

\paragraph{Cross-Attention} 
In~\autoref{tab:crossattn_ablations_1b}, we present ablation studies on the placement of cross-attention modules at different positions within the encoder and decoder of {\textsc BLT}. Specifically, for the encoder cross-attention, we explore three initialization strategies for the queries: (1) assigning the same trainable embedding across all global states, (2) employing a hash embedding generated from the patch's byte sequence, and (3) applying a pooling operation over the encoder's hidden states corresponding to the patch bytes at the relevant layer.

Our results indicate that incorporating cross-attention within the \textit{decoder} yields the strongest performance gains. In contrast, integrating cross-attention into the encoder yields marginal benefits, and only when the queries are initialized using pooling. Furthermore, our analysis shows that cross-attention provides notable advantages when working with Common-Crawl data, with the improvements being most pronounced for larger patch sizes.

\paragraph{n-gram Hash Embeddings} 

We evaluate the impact of varying n-gram hash embedding vocabulary sizes at 0, 100K, 200K, and 400K, with corresponding results shown in Table~\ref{tab:ngram_ablations_1b}. Our observations indicate that the use of hash embeddings benefits all evaluated domains, with especially pronounced gains on Wikipedia and Github (yielding a 0.04 bpb improvement, versus a 0.01 bpb gain observed after 15k steps at the 8B scale). At the 8B parameter setting, adjusting the number of hashes from 500K down to 300K resulted in an approximate performance variation of 0.001 bpb after 15k training steps. These findings suggest that hash embeddings are essential for {\textsc BLT} to achieve parity with tokenizer-based approaches. Nonetheless, increases beyond 300K hashes yield progressively smaller benefits. Furthermore, our results suggest that the enhancements from hash embeddings are mostly orthogonal to those provided by cross-attention, as each method improves performance on distinct subsets of datasets.

\paragraph{Local Model Hyperparamaters} Table \ref{tab:local_layers} presents an ablation study exploring different choices for the number of layers in the local encoder and decoder. Our findings indicate that, in conjunction with hash n-gram embeddings, {\textsc BLT} achieves effective performance when employing a very minimal encoder—specifically, a single-layer encoder—alongside a more substantial, multi-layer decoder.

\section{Related Work}
\label{section:related_work}

\textbf{Character-Level RNNs:} Character Language Modeling has been a key area of research since the introduction of neural network architectures~\citep{sutskever2011generating,mikolov2012subword,graves2013generating}, primarily because it naturally captures out-of-vocabulary words without requiring explicit back-off strategies. \cite{kim2016character} developed a model that exclusively encodes input text at the character level by utilizing a combination of convolutional and highway networks, which subsequently interface with LSTM-based RNNs. Their approach achieved comparable results to then-current RNN-based state-of-the-art language models on English, and even exceeded their performance in morphologically complex languages—highlighting another important benefit of character-level LLMs. In a similar vein, \cite{kenter2018byte} employed byte-level LSTM models for machine comprehension tasks, demonstrating superior results over word-level models for languages with rich morphological structures, such as Turkish and Russian. Likewise, \cite{zhang2015character} explored the use of character-based convolutional networks for classification problems, finding that these models outperformed their word-based counterparts for certain cases. Hierarchical modeling at the character level was further advanced by \citet{chung2019hierarchical}, who leveraged LSTM architectures with hierarchical boundary detection to uncover latent text structures and boost language modeling efficacy. Additionally, ByteNet~\citep{kalchbrenner2016neural} introduced convolutional layers operating over characters, rather than employing attention mechanisms, within the context of machine translation.

\textbf{Character-Level Transformers:} The introduction of attention-based transformer architectures~\citep{vaswani2017attention}, alongside techniques such as subword tokenization~\citep{sennrich-etal-2016-neural}, led to marked advances in language modeling and various standardized tasks. Despite these gains, the use of word and sub-word segments inherently embeds assumptions about the level of abstraction the model should target. In an attempt to integrate the achievements of transformer architectures with encouraging early findings from character-level modeling, \citet{al2019character} constructed very deep transformer models. By employing auxiliary objectives during training, they managed to surpass LSTM-based character language models in performance. Nevertheless, a notable disparity persisted between character-level and word-level large language models. Additionally, GPT-2~\citep{radfordgpt2} found that, even on sizable datasets such as the 1 billion word benchmark, language models operating at the byte level lagged behind their word-level counterparts.

Previous work by \cite{Choe2019BridgingTG} revealed that transformer-based LLMs at the byte level can surpass the performance of subword-level LLMs given equivalent parameter counts. However, their findings also indicated a significantly greater computational cost and much longer training durations for byte-level models. Likewise, \cite{el2020characterbert} introduced CharFormer, a BERT variant constructing word representations through convolutions over character embeddings, achieving notable gains in the biomedical context, albeit at a much higher computational expense. In another effort, \cite{clark2022canine} presented CANINE, a 150M parameter encoder that processes raw character inputs. CANINE features a deep transformer architecture that closely resembles our global model’s structure, and also incorporates a mix of local transformers and strided convolutional layers to efficiently downsample the character input. This model demonstrated improved performance over the corresponding token-level encoder mBERT across various multilingual benchmarks. ByT5~\citep{xue2022byt5} examined alternatives for constructing byte-level encoder-decoder models, notably omitting any patching mechanisms. Their model achieved both improved noise robustness and competitiveness against tokenizer-based architectures using only a quarter of the data. Nevertheless, the absence of patching forced the model to perform computationally intensive attention operations on every byte, resulting in substantial computational demands. Shifting to bytes as modeling units extends sequence lengths, making the efficient scaling of byte-level approaches inherently difficult. Most recently, leveraging the Mamba Architecture~\citep{gu2023mamba}, which is capable of sustaining fixed-size memory across very extensive contexts, \cite{wang2024mambabyte} devised a patchless, byte-level Mamba model. Their approach surpassed byte-level transformer architectures under controlled computation budgets at the 350M parameter scale, achieving superior bits-per-byte results on several datasets.

\textbf{Patching-based approaches:} Leveraging patching strategies can reduce the otherwise high computational cost (in terms of flops) associated with byte-level LLMs, all while maintaining, or even enhancing, performance. Initial studies have yielded promising outcomes for smaller models and limited training byte budgets. For instance, \cite{nawrot-etal-2022-hierarchical} implemented static patching, incorporating downsampling and upsampling techniques, and introduced the hourglass transformer, which surpassed competing byte-level models at the 150M parameter level. Building on this, \cite{nawrot-etal-2023-efficient} introduced dynamic patching mechanisms. These included an end-to-end trainable boundary predictor, a variant of the boundary predictor that is guided by specific tokenizers, and an entropy-driven patching model analogous to {\textsc BLT}. Their work demonstrated that, for language modeling at a 40M parameter scale and training on approximately 400M tokens, these methods could exceed the performance of standard transformer models of the era. In another study, \cite{lester2024training} explored compressing input sequences via arithmetic coding, attaining levels of compression exceeding those made possible by BPE. Through the use of an equal-info windows methodology, they reported results surpassing byte-level benchmarks for language modeling, though performance still lagged behind models utilizing subword units.

Our research is primarily influenced by and most comparable to the MegaByte approach~\citep{yu2023megabyte}, a decoder-only causal language model that implements a predetermined, static method for grouping and merging representations, thereby translating bytes into fixed-size patches. On the decoder side, a local submodel reconverts these patches back to individual bytes. MegaByte's creators demonstrate its ability to achieve performance parity with token-based models, using around 1 billion parameters on data comprising roughly 400B bytes. In our evaluations, we systematically assess MegaByte as a baseline, observing that its static patching mechanism falls short of the state-of-the-art tokenizer-based models optimally trained with equivalent computational resources. Our findings further reveal how {\textsc BLT} can close this performance divide. Similarly, \cite{slagle2024spacebyte} comment on these limitations of MegaByte, proposing enhancements such as incorporating patching over whitespace and analogous characters and supplementing with a local encoder module. Their work shows that, under constrained computational budgets, their modifications outperform tokenized transformer models within certain specialized domains like Github and arXiv at the 1B parameter level. We replicate these experiments and conclude that further advances in model architecture are essential for byte-level models to reach the performance levels of leading token-based systems, exemplified by Llama 3.

\section{Limitations and Future Work}

For our architectural decisions, we trained models for the number of steps deemed optimal for Llama~3~\citep{dubey2024llama}. Nonetheless, these scaling laws are derived from BPE-level transformers and may not yield ideal ratios of data to parameters when applied to {\textsc BLT}. Deriving scaling laws specifically tailored to {\textsc BLT} is left as future work, which may reveal even more advantageous scaling patterns unique to our approach. Moreover, the majority of our experiments were performed at model sizes up to 1B parameters, raising the possibility that optimal architectural settings could shift when scaling to models with 8B parameters or more, potentially resulting in additional performance gains at such larger scales.

Current transformer libraries and codebases have been optimized primarily for tokenizer-based transformer models. Although we provide theoretical comparisons based on matched FLOPs and utilize some optimized components like FlexAttention for non-standard transformer layers, our implementations may still fall short of tokenizer-based models with respect to actual wall-clock performance. Additional optimization work could further improve runtime efficiency.

Although {\textsc BLT} employs an independently trained entropy model for patching, integrating patching model training within an end-to-end learning framework represents a promising avenue for future research. In Section \ref{sec:distillation}, we describe preliminary studies that demonstrate potential for converting tokenizer-based models like Llama 3 (trained on datasets exceeding 10T tokens) to byte-based models. This is achieved by transferring their weights to the global transformer and keeping them fixed during initialization. Advancing this line of inquiry could lead to approaches that not only preserve the strengths of bytefying but also surpass the performance of original tokenizer-based architectures—without requiring a complete retraining from scratch.

\section{Conclusion}
\label{section:conclusion}

This work introduces the Byte Latent Transformer (\textbf{BLT}), an innovative framework that challenges the traditional reliance on static vocabulary-based tokenization in large language models. Leveraging a flexible, learnable mechanism to segment bytes into variable-length patches, {\textsc BLT} dynamically tailors computational effort according to the underlying data complexity, thereby yielding notable gains in both computational efficiency and model resilience. Through comprehensive scaling experiments, we establish that {\textsc BLT} achieves performance parity with established tokenization-oriented models such as Llama 3 at model sizes up to 8B and data volumes of 4T bytes, and can attain up to 50\% savings in inference FLOPs with only marginal compromises in evaluation scores. Additionally, {\textsc BLT} introduces a new axis for scaling: users may simultaneously expand model capacity and patch dimension without exceeding a predetermined computational ceiling during inference, a property particularly beneficial in resource-constrained scenarios encountered in practice. By directly processing raw byte sequences, {\textsc BLT} also enhances the system’s capability to manage rare and atypical data, leading to improved robustness against noisy inputs and fostering a more granular comprehension of sub-word information. Collectively, these advances establish {\textsc BLT} as a compelling, scalable alternative to standard tokenization schemes, offering a versatile and robust approach for efficient, adaptive language modeling.

\end{document}